\usepackage[utf8]{inputenc}
%Formato de citación APA: se conserva thebibliography estándar porque el
%contenido usa bibliografía manual sin \cite
\let\stdthebibliography\thebibliography
\let\endstdthebibliography\endthebibliography
\usepackage{apacite}
\let\thebibliography\stdthebibliography
\let\endthebibliography\endstdthebibliography
\bibliographystyle{apacite}
\setcounter{secnumdepth}{3} % Numera los subsubsecciones
\usepackage[T1]{fontenc}
\usepackage[spanish, es-tabla]{babel} %% reemplaza "cuadro" por "tabla"
\decimalpoint  %Cambia coma por punto
\usepackage{mathptmx}
\usepackage{layouts} %Saber ancho de hoja.
\usepackage{fontspec}
\usepackage{amssymb}
\usepackage{graphicx}
\usepackage{changepage} %Agregar identación
\setmainfont{Arial}
\usepackage{setspace}
\usepackage[left=4 cm,right=3 cm,top=3 cm,bottom=3.5 cm]{geometry}
\usepackage{fancyhdr} %Encabezado y Pie de página (1)
\pagestyle{fancy} %Encabezado y Pie de página (2)
\usepackage{lipsum} %Crear texto RAMDOM
\renewcommand{\labelitemi}{$\bullet$} %Circulos para viñetas
\usepackage{titlesec} %Titulos de SECCIONES
\usepackage{tocloft} %Titulos de ÍNDICES
\usepackage[colorlinks=true,linkcolor=negro,citecolor = negro]{hyperref}
\usepackage[mathbf=sym]{unicode-math} % Mantener fuentes matemáticas
\makeatletter  %Comando para REDUCIR ERRORES
\usepackage{multirow} % Agregar TABLAS 
\usepackage{array} % Dar formato a las TABLAS
\usepackage{subcaption} % Insertar SubImagenes
\usepackage{tikz} % Diagrama de Flujo
\usetikzlibrary{calc,positioning,shapes.geometric,shapes.symbols,shapes.misc,arrows.meta,decorations.pathreplacing}

% Insertar formas para diagramas de flujo
\tikzstyle{startstop} = [rectangle, rounded corners, minimum width=3cm, minimum height=0.5cm,text centered, draw=black]
\tikzstyle{io} = [trapezium, trapezium left angle=70, trapezium right angle=110, minimum width=3cm, minimum height=0.5cm, text centered, text width=3cm, draw=black]
\tikzstyle{process} = [rectangle, minimum width=3cm, minimum height=0.5cm, text centered, text width=4cm, draw=black]
\tikzstyle{decision} = [diamond, minimum width=3cm, minimum height=0.5cm, text centered, draw=black]
\tikzstyle{loop} = [chamfered rectangle,chamfered rectangle xsep=2cm,draw=black]
\tikzstyle{arrow} = [->,>=stealth]
\tikzstyle{line}=[draw]

\usepackage{listings}
\usepackage{color}

%New colors defined below
\definecolor{codegreen}{rgb}{0,0.6,0}
\definecolor{codegray}{rgb}{0.5,0.5,0.5}
\definecolor{codepurple}{rgb}{0.2,0,1}
\definecolor{codeRojo}{rgb}{0.7,0,0.3}
\definecolor{backcolour}{rgb}{1.0, 1.0, 1.0}

%Code listing style named "mystyle"
\lstdefinestyle{mystyle}{
  backgroundcolor=\color{backcolour},  
  commentstyle=\color{codegreen},
  keywordstyle=\color{codeRojo},
  numberstyle=\tiny\color{codegray},
  stringstyle=\color{codepurple},
  basicstyle=\footnotesize,
  breakatwhitespace=false,         
  breaklines=true,                 
  captionpos=b,                    
  keepspaces=true,                 
  numbers=left,                    
  numbersep=5pt,                  
  showspaces=false,                
  showstringspaces=false,
  showtabs=false,                  
  tabsize=2
}

%"mystyle" code listing set
\lstset{style=mystyle}

\newenvironment{MyFont}{\fontfamily{ugm}\selectfont}{\par}

\usepackage{changepage} %Agregar espacio a Listing

% --- Paquetes adicionales requeridos por el contenido ---
\usepackage{amsmath}
\usepackage{booktabs}
\usepackage{longtable}
\usepackage{tabularx}
\usepackage{enumitem}
\setlist{topsep=4pt,itemsep=2pt}
\usepackage{pgfplots}
\pgfplotsset{compat=1.18}
\usepackage{pgf-pie}

% Metadatos PDF
\hypersetup{
  pdftitle={Aplicación de una metodología con criterios técnicos y regulatorios para la selección de tecnologías de interfaz de usuario en instituciones bancarias},
  pdfauthor={Willian José Pillaca Meneses},
  pdfsubject={Tesis de Ingeniería de Sistemas},
  pdflang={es-PE}
}

% Datos generales usados por permisos y otros archivos de contenido
\newcommand{\tesisTitulo}{Aplicación de una metodología con criterios técnicos y regulatorios para la selección de tecnologías de interfaz de usuario en instituciones bancarias}
\newcommand{\tesisAutor}{Willian José Pillaca Meneses}
\newcommand{\tesisAnio}{2026}


% Centrado del título del ÍNDICE / LISTA DE FIGURAS / LISTA DE CUADROS

\renewcommand{\cfttoctitlefont}{\hfill \normalfont\normalsize\bfseries}
\renewcommand{\cftaftertoctitle}{\hfill}

\renewcommand{\cftlottitlefont}{\hfill\normalfont\normalsize\bfseries}
\renewcommand{\cftafterlottitle}{\hfill}

\renewcommand{\cftloftitlefont}{\hfill\normalfont\normalsize\bfseries}
\renewcommand{\cftafterloftitle}{\hfill}


% Formato de los CAPÍTULOS, SECCIONES Y SUBSECCIONES
% La clase y fancyhdr no definen \chaptermark, sin esto el encabezado
% quedaría fijado en la última marca (p. ej. LISTA DE FIGURAS)
\renewcommand{\chaptermark}[1]{\markboth{#1}{}}

\titleformat{\chapter}[display]{\normalfont\normalsize\bfseries}{\centering CAPÍTULO \Roman{chapter}}{6pt}{\centering}
\titlespacing*{\chapter}{0pt}{-10pt}{12pt}

\titleformat{\section}[block]{\normalfont\normalsize\bfseries}{\thesection}{0.5em}{}
\titlespacing*{\section}{0pt}{10 pt}{5 pt}

\titleformat{\subsection}[block]{\normalfont\normalsize\bfseries}{\thesubsection}{0.5em}{}
\titlespacing*{\subsection}{0pt}{10 pt}{5 pt}

\titleformat{\subsubsection}[block]{\normalfont\normalsize\bfseries}{\thesubsubsection}{0.5em}{}
\titlespacing*{\subsubsection}{0pt}{10 pt}{5 pt}


% Espaciado entre PÁRAFOS y SANGRÍA
\setlength{\parskip}{5 pt}
\setlength{\parindent}{0cm}


% Dar formato a las TABLAS
\newcolumntype{M}[1]{>{\centering\arraybackslash}m{#1}}
\newcolumntype{L}[1]{>{\raggedright\arraybackslash}m{#1}}
\newcolumntype{R}[1]{>{\raggedleft\arraybackslash}m{#1}}
\newcolumntype{N}{@{}m{0pt}@{}}
\renewcommand{\arraystretch}{1.25}

% Cambiar titulo de bibliografía
\addto\captionsspanish{\renewcommand{\bibname}{\centering REFERENCIAS BIBLIOGRÁFICAS}}


% Posicionamiento vertical de TOC, LOT and LOF

\setlength{\cftbeforelottitleskip}{1pt} 
\renewcommand{\cftafterlottitleskip}{12pt}

\setlength{\cftbeforeloftitleskip}{1pt} 
\renewcommand{\cftafterloftitleskip}{12pt}

\setlength{\cftbeforetoctitleskip}{16pt} 
\renewcommand{\cftaftertoctitleskip}{12 pt}


% Cambiando las etiquetas de las FIGURAS y TABLAS (Caption y autoref)
\addto\captionsspanish{\renewcommand{\figurename}{\footnotesize FIGURA N°}}
\addto\extrasspanish{\def\figureautorefname{ Figura N°}}
 
\addto\captionsspanish{\renewcommand{\tablename}{\footnotesize TABLA N°}}
\addto\extrasspanish{\def\tableautorefname{ Tabla N°}} 


% Espaciamiento dentro del índice

\setlength{\cftbeforechapskip}{2mm}
\renewcommand\cftchapafterpnum{\vskip6pt}
\renewcommand\cftsecafterpnum{\vskip5pt}
\renewcommand\cftsubsecafterpnum{\vskip5pt}


%Agregar la palabra CAPITULO al TOC, FIGURA a LOF y TABLA al LOT

\renewcommand{\cftchappresnum}{CAPÍTULO }
\renewcommand{\cftchapaftersnum}{:}
\renewcommand{\cftchapnumwidth}{7em}

\renewcommand{\cftfigpresnum}{Figura N° }
%\renewcommand{\cftfigaftersnum}{:}
\renewcommand{\cftfignumwidth}{6.85 em}

\renewcommand{\cfttabpresnum}{Tabla N° }
%\renewcommand{\cftfigaftersnum}{:}
\renewcommand{\cfttabnumwidth}{6.5 em}


%Definición de COLORES

\usepackage{xcolor}
\definecolor{granate}{RGB}{113,22,16}
\definecolor{gris}{RGB}{154,153,157}
\definecolor{arena}{RGB}{230,217,170}
\definecolor{azul}{rgb}{0.03, 0.15, 0.4}
\definecolor{negro}{rgb}{0, 0, 0}

%Cambiando a Números Romanos los Capítulos
\renewcommand{\thechapter}{\Roman{chapter}}
\renewcommand{\theequation}{\arabic{chapter}.\arabic{equation}} 
\renewcommand{\thesection}{\arabic{chapter}.\arabic{section}}  
\renewcommand{\thetable}{\arabic{chapter}.\arabic{table}}  
\renewcommand{\thefigure}{\arabic{chapter}.\arabic{figure}}



%Escribir los INFORMACIÓN PERSONAL Y DEL TRABAJO

%Autor para PIE DE PÁGINA (Respetar mayusculas y minisculas)
\author{WILLIAN JOSÉ PILLACA MENESES}

%Autor para CARÁTULA (Siempre en mayuscula y sin saltos de linea)
\authorcaratula{WILLIAN JOSÉ PILLACA MENESES}

%Título en para el PIE DE PÁGINA (versión corta: el título completo desborda
%el pie de página y colisiona con el número de página)
\title{Selección de tecnologías de interfaz de usuario en instituciones bancarias}

%Título para CARÁTULA (Siempre en mayuscula y sin saltos de linea)
\titlecaratula{APLICACIÓN DE UNA METODOLOGÍA CON CRITERIOS TÉCNICOS Y REGULATORIOS PARA LA SELECCIÓN DE TECNOLOGÍAS DE INTERFAZ DE USUARIO EN INSTITUCIONES BANCARIAS}

%Nombre de la FACULTAD (Siempre en mayuscula)
\facultad{FACULTAD DE INGENIERÍA INDUSTRIAL Y DE SISTEMAS}

%Para obtener el título profesional de ... (Siempre en mayuscula)
\grado{INGENIERO DE SISTEMAS}

%Asesor para CARÁTULA (Siempre en mayuscula y sin saltos de linea)
\asesor{DR. NOMBRE DEL ASESOR}

%AÑO para CARÁTULA
\yyearr{2026}
